%===============================================================
% CHAPTER 7 – Preliminary Design Review (PDR) Document
%===============================================================
\chapter{Preliminary Design Review (PDR) Document}

\section{Design Concept Summary}

At the time of this Preliminary Design Review, the Solift H-VTOL design is at Technology Readiness Level (TRL) 3--4. The primary design concept -- a three-axis solar-powered tilt-rotor with a 12S heterogeneous drive train and ETFE-encapsulated IBC solar cells on the upper wing surface -- has been validated through:

\begin{itemize}[noitemsep]
    \item Phase~1 material selection matrix (Section~4.1), confirming suitability of CFRP/Nomex composites, aluminium alloys, and ETFE encapsulation for manufacturability and durability.
    \item Phase~2 full parametric SolidWorks model with confirmed geometric envelope compliance, ensuring payload bay dimensions, wing span constraints, and CG envelope are satisfied.
    \item Phase~3 FEA confirming $\text{SF} = 2.69$ at the critical tilt-arm root (Table~\ref{tab:fea}), demonstrating adequate structural safety margin under combined bending and torsion loads.
    \item Phase~3 CFD confirming $L/D = 45.6$ at cruise (Table~\ref{tab:cfd}), validating aerodynamic efficiency and range potential under the selected MH114 airfoil.
    \item Phase~4 tolerance stack-up confirming $\delta_{total} = 0.60^\circ < 1.0^\circ$ requirement, ensuring tilt mechanism precision and vibration control during transition.
    \item Phase~5 manufacturing process plan with $\sim$39~h estimated production time per prototype unit, including CNC machining, composite layup, and subsystem assembly.
\end{itemize}

\noindent The design is recommended to proceed to \textbf{Critical Design Review (CDR)} subject to the action items in Section~\ref{sec:pdr_actions}.

\section{Key Trade-Off Studies}

Three key trade-off decisions were analysed and resolved during the design phase:

\subsection{Trade-Off 1: Battery Mass vs.\ Solar Panel Area vs.\ Payload Capacity}

The range equation for the Solift in cruise can be written as:
\begin{equation}
    R = \frac{\eta_{prop} \cdot (E_{battery} + E_{solar})}{m_{MTOW} \cdot g} \cdot \frac{L}{D}
    \label{eq:range}
\end{equation}

A parametric sweep (Figure~\ref{fig:tradeoff_range}) shows:
\begin{itemize}[noitemsep]
    \item Increasing battery from 22~Ah to 30~Ah adds $\sim$18~km range but reduces payload from 5.0 to 3.2~kg due to higher mass.
    \item Increasing solar area from 0.95 to 1.3~m$^2$ (requiring a 20\% wingspan increase) adds $\sim$22~km range while keeping payload at 5.0~kg, but adds 0.8~kg structural mass and complicates wing manufacturing.
    \item \textbf{Decision:} Retain 22~Ah battery and 0.95~m$^2$ solar area as the best balance for the 120~km / 5~kg baseline requirement, ensuring manufacturability and mission compliance.
\end{itemize}

\subsection{Trade-Off 2: CNC vs.\ 3D-Printed Tilt Mechanism}

\begin{table}[H]
\centering
\caption{Manufacturing Trade-Off: Tilt Mechanism Fabrication Method}
\label{tab:tradeoff_mfg}
\begin{tabularx}{\textwidth}{L{4.5cm} C{3cm} C{3cm} C{3cm}}
\toprule
\textbf{Criterion} & \textbf{CNC -- Al 7075-T6} & \textbf{DMLS -- Ti64} & \textbf{FDM -- CF-Nylon} \\
\midrule
Mass per arm (g) & 142 & 118 & 98 \\
Fatigue life (estimated cycles) & $> 10^7$ & $> 10^7$ & $< 10^5$ \\
Cost per arm (USD, prototype) & 280 & 950 & 35 \\
Lead time (days) & 3 & 14 & 0.5 \\
Dimensional accuracy & $\pm$0.05~mm & $\pm$0.1~mm & $\pm$0.3~mm \\
\textbf{Decision} & \textbf{Selected} & Rejected (cost) & Rejected (fatigue) \\
\bottomrule
\end{tabularx}
\end{table}

\noindent CNC machining of 7075-T6 aluminium was selected for its balance of strength, fatigue life, and cost. DMLS titanium offered weight savings but was prohibitively expensive, while FDM CF-Nylon was cost-effective but failed fatigue requirements.

\subsection{Trade-Off 3: Airfoil Selection}

Three airfoils were evaluated at $Re = 3 \times 10^5$:
\begin{itemize}[noitemsep]
    \item \textbf{MH114:} $C_{L,max} = 1.52$, $L/D_{max} = 46.3$. \textbf{Selected for its high lift-to-drag ratio, benign stall behaviour, and compatibility with conformal solar cell integration.}
    \item Clark-Y: $C_{L,max} = 1.41$, $L/D_{max} = 41.0$. Inferior $L/D$; rejected.
    \item Selig S1223: $C_{L,max} = 2.21$, $L/D_{max} = 38.5$. High lift but low $L/D$; excessive curvature unsuitable for solar cell integration.
\end{itemize}

\section{Risk Assessment Matrix}

\begin{table}[htbp]
\centering
\caption{PDR Risk Assessment Matrix}
\label{tab:risk}
\small
\begin{tabularx}{\textwidth}{C{1cm} L{3.5cm} C{1.4cm} C{1.4cm} C{1cm} X}
\toprule
\textbf{ID} & \textbf{Risk Description} & \textbf{Likelihood (1--5)} & \textbf{Severity (1--5)} & \textbf{RPN} & \textbf{Mitigation} \\
\midrule
R-01 & VTOL-to-cruise power failure causing stall during transition & 2 & 5 & 10 & Dual-motor redundancy; abort-to-hover logic; geofenced flight test zone \\
R-02 & CNC tilt-arm precision defect causing excess vibration & 2 & 4 & 8 & CMM inspection of all arms; vibration spectrum test pre-flight \\
R-03 & Solar cell delamination under thermal cycling & 3 & 3 & 9 & ETFE vacuum lamination; ground-truth temperature cycling test (50 cycles, $-10$ to $+55^\circ$C) \\
R-04 & Battery thermal runaway during fast recharge & 2 & 5 & 10 & BMS with cell-level temperature monitoring; C-rate limited to 1.5C \\
R-05 & CFRP wing delamination at spar root & 1 & 5 & 5 & FEA-verified joint design; proof-load test to 1.5$\times$ MTOW \\
R-06 & Flight controller firmware instability during tilt transition & 3 & 4 & 12 & Hardware-in-the-loop (HIL) simulation; Pixhawk 6X with ArduPlane VTOL mode \\
\bottomrule
\multicolumn{6}{l}{\small RPN = Risk Priority Number = Likelihood $\times$ Severity} \\
\end{tabularx}
\end{table}

\section{Review Decisions and Action Items}
\label{sec:pdr_actions}

The following action items were raised during the PDR and must be resolved prior to CDR:

\begin{table}[H]
\centering
\caption{PDR Action Items}
\label{tab:actions}
\begin{tabularx}{\textwidth}{C{1cm} L{5cm} L{3cm} C{2cm} C{1.5cm}}
\toprule
\textbf{AI\#} & \textbf{Action Required} & \textbf{Owner} & \textbf{Due Date} & \textbf{Status} \\
\midrule
AI-01 & Complete 3D FEA of full tilt arm under combined bending + torsion & [Name 1] & [Date] & Open \\
AI-02 & Procure Sunnysky X4112S motors and measure thrust curve at 12S & [Name 2] & [Date] & Open \\
AI-03 & Perform thermal cycling test on ETFE/IBC coupon (50 cycles) & [Name 3] & [Date] & Open \\
AI-04 & Validate VTOL-to-cruise transition in HIL simulation (ArduPlane) & [Name 4] & [Date] & Open \\
AI-05 & Update SRR traceability matrix to reflect Phase~4 tolerance decisions & [Name 1] & [Date] & Open \\
AI-06 & Obtain supplier quotations for 7075-T6 billet and CNC machining & [Name 2] & [Date] & Open \\
\bottomrule
\end{tabularx}
\end{table}

\noindent \textbf{Overall PDR Verdict:} The Solift H-VTOL design is approved to proceed to the fabrication phase, pending closure of AI-01, AI-02, and AI-04 (flight-critical items). These action items directly affect structural integrity, propulsion performance, and transition stability, and therefore must be resolved before Critical Design Review. All remaining action items (AI-03, AI-05, AI-06) are considered secondary but will be reviewed at CDR to ensure manufacturability, environmental durability, and supply chain readiness.

\noindent The PDR confirmed that:
\begin{itemize}[noitemsep]
    \item The design concept meets baseline mission requirements (5~kg payload, 120~km range, VTOL capability).
    \item Structural and aerodynamic analyses provide sufficient safety margins for tilt-arm loads and cruise efficiency.
    \item Manufacturing trade-offs (CNC vs.\ additive vs.\ polymer) were resolved in favor of CNC aluminium for reliability and cost balance.
    \item Risk mitigation strategies are in place for critical subsystems (tilt mechanism, solar integration, battery safety).
    \item Verification methods (flight tests, FEA/CFD, tolerance stack-up, environmental cycling) are aligned with SRR requirements.
\end{itemize}

\noindent \textbf{Next Steps:}  
The team will now focus on closing open action items, finalizing supplier contracts, and preparing prototype fabrication. The Critical Design Review (CDR) will emphasize:
\begin{enumerate}[noitemsep]
    \item Validation of tilt-arm FEA under combined load cases.
    \item Empirical thrust curves for selected motors at 12S voltage.
    \item Hardware-in-the-loop validation of VTOL-to-cruise transition logic.
    \item Confirmation of ETFE/IBC solar coupon durability under thermal cycling.
    \item Updated traceability matrix reflecting tolerance and manufacturability decisions.
\end{enumerate}

\noindent With these steps completed, Solift will advance from TRL~3--4 toward TRL~5, entering the prototype fabrication and subsystem integration stage.
